\section{Introduction}
\label{sec:introduction}

The Standard Model and general relativity are extraordinarily successful effective descriptions of observed physics.
At the same time, their combined foundational inventory---multiple quantum fields, symmetry sectors, and parameters fixed by experiment---continues to motivate a complementary question:
could some of that effective complexity arise from a smaller microscopic set of ingredients?
The present paper takes up that question in a deliberately limited form.

\paragraph{Substrate-first strategy.}
We adopt a deterministic \emph{substrate-first} viewpoint.
The primitive ontology is a single dynamical medium, and familiar field-theoretic objects are interpreted as coarse-grained descriptions of special excitation sectors of that medium.
The aim is not to replace established theories by assertion.
The aim is to formulate a mathematically coherent theory fragment in which relativistic kinematics, geometric gauge variables, and particle-like localization can all be discussed within one substrate model.

\paragraph{Minimal commitments.}
The framework is built from four basic commitments:
(i) a three-dimensional brane-like medium with local action dynamics,
(ii) a 4D embedding so that transverse excitation can modify intrinsic geometry,
(iii) a cubic lattice as the candidate microscopic realization, and
(iv) an ordered narrowband sector whose transported polarization states can carry geometric phase.
Everything beyond that---effective gauge potentials, charge-like labels, Lorentz symmetry, and particle-like spectra---is treated as emergent bookkeeping rather than as additional fundamental substances.

\paragraph{Continuum/lattice split.}
A recurring source of confusion in speculative substrate models is the mixing of continuum and microscopic claims.
Here the distinction is made explicit.
The continuum theory is used to state the geometric and constitutive content in a coordinate-free way.
The cubic lattice is the concrete ultraviolet hypothesis that determines branch splitting, anisotropy, Brillouin-zone structure, and the axis-aligned mixing mechanisms that the continuum description cannot by itself resolve.
This paper therefore treats the continuum equations as the long-wavelength effective theory of a lattice substrate, not as a replacement for it.

\paragraph{Why the localization problem matters.}
The decisive issue for any such framework is not whether one can write down a plausible continuum action.
It is whether the same substrate can support robust localized modes with the right scales, internal labels, and long-range diagnostics.
In earlier exploratory work, electron-like localization seemed the natural first target.
At the present stage, however, the more informative target is a proton- or neutron-like state:
the cubic substrate naturally offers a three-channel axis-aligned sector, so baryon-scale triplet states provide the clearest setting in which branch mixing, confinement, and non-Abelian holonomy can be tested together.

\paragraph{Relation to existing ideas.}
The emergence of an effective light cone from a preferred-time medium has close analogues in analogue-gravity systems \citep{Barcelo2011AnalogueGravity,Volovik2003HeliumDroplet,Jacobson2004EinsteinAether}.
The use of Berry and Wilczek--Zee connections follows standard geometric-phase theory in wave systems \citep{Pancharatnam1956,Berry1984,WilczekZee1984,Bliokh2015SpinOrbit,Cisowski2022}.
The present use of the term ``brane'' is literal---an elastic medium embedded in $\R^4$---and should not be confused with high-energy brane-world scenarios \citep{RandallSundrum1999}.

\paragraph{Paper organization.}
\Cref{sec:scope} states the precise scope, core claims, and non-claims.
\Cref{sec:continuum-model} formulates the continuum substrate model and its lattice interpretation.
\Cref{sec:emergent-relativity} develops the linear-wave regime and its effective Lorentz interpretation.
\Cref{sec:geophase,sec:effective-field} isolate the geometric-phase sector and the corresponding envelope description.
\Cref{sec:localized} formulates localized particle-like states as nonlinear bound modes.
\Cref{sec:computational-program} turns these ingredients into a concrete numerical program, with baryon-scale triplet states as the immediate target.
\Cref{sec:discussion} and \Cref{sec:conclusion} summarize what is actually established and what remains open.
